\section{Colour: logistic-map instability}\label{sec:colour}
\stagebox{\statusplanned}{none yet}

Colour is not part of the Markov state; it is a separate future stage for buds
with several differently coloured or colour-graded petals. This section
records the design.

\subsection{The logistic map}
\[
  x_{n+1} = r\,x_n(1 - x_n), \qquad x_n\in[0,1],\; r\in[0,4].
\]
For $1 < r < 3$ every orbit settles on the fixed point $x^\star = 1 - 1/r$. At
$r_1 = 3$ it splits into a 2-cycle, at $r_2 = 1+\sqrt6 \approx 3.449$ into a
4-cycle, and the period-doublings accumulate at $r_\infty\approx 3.5699$, beyond
which most $r$ give chaos. How fast nearby orbits separate is measured by the
Lyapunov exponent
\[
  \Lambda(r) = \lim_{N\to\infty}\frac1N\sum_{n=0}^{N-1}\ln\lvert r(1 - 2x_n)\rvert,
\]
negative in stable regimes and positive in chaotic ones.

\subsection{How the model will use it}
\begin{itemize}[leftmargin=*]
\item Each petal, or each region of a petal, inherits a position $r$ on the
  parameter axis.
\item Far inside a stable regime, small random differences die out, so the
  petal gets the same colour in every flower on the plant.
\item Near a bifurcation, or in the chaotic band, small stochastic differences
  are amplified, so only those petals vary between sibling flowers.
\item The Markov chain does not wander into instability by itself. Its steps
  drive the logistic iterations from the inherited point, so unstable petals
  look chaotic and stable ones stay locked.
\end{itemize}

\subsection{Inheritance (later)}
Part of what a plant passes on is where its petals sit on the bifurcation
diagram, which makes ``closeness to chaos'' a heritable, selectable trait. This
waits until the single-plant model works.
